\section{Chapitre~\ref{sec:recognition}. Reconnaissance}

La vitesse de reconnaissance, la structure des premiers étages, la lecture linéaire de la réponse et l'apprentissage sur la continuité du temps sont mesurés chez l'homme et le singe; la réduction de toute la chaîne à deux opérations et l'apprentissage sur la vidéo sont vérifiés sur les modèles du livre; les explications des illusions vont du bien fondé au discuté.

{\footnotesize
\setlength{\tabcolsep}{3pt}\renewcommand{\arraystretch}{1.15}
\begin{longtable}{>{\raggedright}p{0.5cm}>{\raggedright}p{4.0cm}>{\raggedright}p{5.6cm}>{\raggedright}p{2.3cm}>{\raggedright\arraybackslash}p{1.7cm}}
\toprule
N\textsuperscript{o} & Affirmation & Expérience et ce qui est mesuré & Source & Statut \\
\midrule\endfirsthead
\toprule
N\textsuperscript{o} & Affirmation & Expérience et ce qui est mesuré & Source & Statut \\
\midrule\endhead
1 & Sous translation, la comparaison pixel par pixel avec un modèle ne fait pas mieux que pile ou face; la paire d'étages~\eqref{eq:scstage} reconnaît la figure à n'importe quelle position & \texttt{models/\allowbreak recognition\_models.py}, \og rétine\fg{} de $24$ points, $4000$ essais: modèle, $0.49$, $0.51$, $0.52$ pour une proportion de points inversés de $0$, $0.1$, $0.2$; paire $S$ et $C$, $1.00$, $0.86$, $0.70$ & déduction du chapitre & modèle \\
2 & Un animal sur une image est reconnu en $\sim150$\,ms, en un seul passage à travers une dizaine d'étages de $10$--$20$\,ms & Homme, tâche \og animal ou pas\fg{} sur des photos présentées $20$\,ms: les potentiels évoqués des deux réponses divergent après $\sim150$\,ms. Singe: le temps de première réponse croît d'étage en étage (V1 d'abord, puis V2 et V4). Le nombre d'étages et \og zéro ou une impulsion par étage\fg{} sont déduits de ces chiffres & \cite{Thorpe1996,Schmolesky1998} & mesuré \\
3 & L'étage $S$ est un ET sur les parties, l'étage $C$ un maximum sur les positions (proposition~\ref{prop:conveyor}) & Chat, cortex visuel primaire: les cellules simples répondent à un bord d'orientation donnée à un endroit donné, les cellules complexes à la même orientation sur une zone plus grande. Enregistrement intracellulaire de cellules complexes: la réponse à deux barres est, pour beaucoup de cellules, proche de la plus grande des deux réponses, en moyenne très proche de l'opération max. Maximum par inhibition mutuelle: proposition~\ref{prop:wta}; division par l'activité totale: revue & \cite{HubelWiesel1962,Lampl2004,Carandini2012} & mesuré \\
4 & Plus haut dans la chaîne, les combinaisons sont plus grandes et plus complexes, et la réponse dépend de moins en moins de la position & Singe, simplification des stimuli jusqu'au \og trait critique\fg{}: les cellules de V4 et du cortex inférotemporal postérieur répondent à des combinaisons complexes de forme, de couleur et de texture avec un petit champ; dans le cortex inférotemporal antérieur, le champ est plus grand. L'alternance de $S$ et $C$ dans le modèle reproduit ce tableau & \cite{KobatakeTanaka1994,Riesenhuber1999} & mesuré \\
5 & Les réseaux convolutifs reproduisent cette alternance et prédisent le mieux les réponses des étages supérieurs; l'objet se lit linéairement à partir des fréquences d'un groupe de neurones & Réseau entraîné à reconnaître, sans ajustement au cerveau: la couche supérieure prédit les réponses des neurones du cortex inférotemporal du singe, les couches intermédiaires celles de V4. Un classifieur linéaire sur l'activité de $\sim100$ cellules prises au hasard, dans des fenêtres dès $12.5$\,ms, reconnaît l'objet et sa catégorie à différentes positions et tailles & \cite{Fukushima1980,Yamins2014,Hung2005} & mesuré \\
6 & La règle de Hebb avec trace~\eqref{eq:traceinvariance} apprend l'invariance sur une vidéo sans étiquettes si la trace dure au moins $\sim20$\,ms & L'idée des traits lents et de la règle avec trace relève de la théorie. \texttt{models/\allowbreak recognition\_models.py}, deux neurones $C$, $16$ entrées, $10$ essais, pause entre objets de $0.3$\,s. Pour $\tau_{\text{tr}}=5$\,ms, coïncidence avec la figure de $0.52$ en moyenne, pour $10$\,ms, $0.56$; pour $20$, $50$ et $200$\,ms, $1.00$ dans tous les essais (à $30$\,ms, un essai sur dix se trompe) & \cite{Foldiak1991,WiskottSejnowski2002} & modèle \\
7 & Dans le cerveau, l'invariance s'apprend sur la continuité du temps & Singe: on substituait discrètement l'objet pendant un saut de l'œil vers un endroit donné du champ; l'invariance de position des neurones du cortex inférotemporal se modifiait conformément à la substitution, de façon significative dès une heure. Qu'il s'agisse de la règle principale de l'apprentissage visuel est une hypothèse & \cite{LiDiCarlo2008} & mesuré \\
8 & Mouvement: détecteurs de direction, puis la même paire ET/OU sur de grandes zones & Détecteurs de direction dans la rétine du lapin; dans l'aire MT du singe, les $110$ neurones enregistrés sont tous sélectifs à la direction, et la réponse au mouvement y est plus forte que dans V1 & \cite{Barlow1965,Albright1984} & mesuré \\
9 & Les étages supérieurs envoient une prédiction vers le bas, l'erreur monte & Le modèle explique les effets extra-classiques des champs du cortex visuel primaire; certaines observations concordent (revue), mais ce n'est pas démontré comme principe général du pipeline & \cite{RaoBallard1999,Keller2018} & hypothèse \\
10 & Les images difficiles exigent des rétroactions et plusieurs tours & Singe: pour les images \og difficiles\fg{}, la décision dans le cortex inférotemporal se forme $\sim30$\,ms plus tard; ces réponses tardives sont mal prédites par un réseau direct, mieux par des réseaux récurrents ou plus profonds & \cite{Kar2019} & mesuré \\
11 & Une falsification invisible des pixels trompe le réseau; la vision humaine est bien plus robuste, mais pas invulnérable & Sur les réseaux: une petite perturbation choisie tout exprès change la réponse; l'exemple \og panda $\to$ gibbon\fg{} vient du deuxième article. Chez l'homme, en présentation brève, des falsifications qui se transfèrent bien d'un réseau à l'autre infléchissent le choix & \cite{Szegedy2014,Goodfellow2015,Elsayed2018} & mesuré \\
12 & Illusions d'attente: une entrée peu fiable cède à l'attente; ce qui est pâle bouge \og plus lentement\fg{}, la robe est vue de façons différentes (section~\ref{sec:illusions}) & Des réseaux de barres moins contrastés semblent bouger plus lentement. Enquête auprès de $1401$ personnes: la robe est vue selon trois variantes stables, un indice sur l'éclairage fait basculer la couleur; chez $\sim13\,000$ observateurs, c'est l'hypothèse sur l'éclairage qui décide. Un modèle bayésien avec attente de mouvements lents explique plusieurs illusions de mouvement & \cite{Thompson1982,LaferSousa2015,Wallisch2017,Weiss2002,Purves2011} & hypothèse \\
13 & Réglage du gain: cascade, inclinaison et contraste; la grille d'Hermann ne s'explique pas par l'inhibition des voisins & Les détecteurs de direction de la rétine du lapin, lors d'un mouvement prolongé, baissent leur fréquence et, après l'arrêt, descendent sous le niveau de fond. L'illusion d'inclinaison est reproduite par un modèle de division par l'activité des voisins. Des modifications de la grille qui ne changent pas la réponse des neurones de sortie de la rétine suppriment les taches: l'explication par inhibition des voisins est rejetée & \cite{BarlowHill1963,Schwartz2009,SchillerCarvey2005} & mesuré \\
14 & Compensation des délais: un objet en mouvement semble en avance sur l'éclair & L'illusion est mesurée. La psychophysique contredit à la fois le prolongement du mouvement et la différence de délais; une explication a posteriori est proposée, par les événements survenant $\sim80$\,ms après l'éclair. Le débat n'est pas tranché & \cite{Nijhawan1994,EaglemanSejnowski2000} & hypothèse \\
15 & Des \og serpents\fg{} immobiles tournent: la différence de délais~\eqref{eq:latency} est lue par les détecteurs de direction & L'illusion fonctionne aussi sans couleur; des paires d'éléments voisins la produisent isolément; des neurones corticaux du singe sélectifs à la direction donnent une réponse orientée à ces paires immobiles. Chez l'homme, la rotation est déclenchée par les microsaccades et les clignements & \cite{Conway2005,OteroMillan2012} & mesuré \\
16 & Images ambiguës: inhibition mutuelle, fatigue et bruit; les bascules sont surtout dues au bruit & Modèle de groupes de neurones en compétition: les bascules y sont causées par le \emph{bruit}, sans bruit il n'y en a pas; il explique la hausse de la fréquence des bascules avec l'intensité du stimulus. La fatigue y joue un rôle moindre & \cite{MorenoBote2007} & hypothèse \\
17 & Une partie des illusions découle de la tâche qu'apprend la machine & Un réseau entraîné à prédire l'image suivante d'une vidéo \og voit\fg{} la rotation d'anneaux immobiles et ne la voit pas sur les images témoins; des réseaux de débruitage et d'augmentation de la netteté cèdent aux illusions de couleur et de luminosité, comme l'homme & \cite{Watanabe2018,GomezVilla2020} & mesuré \\
\bottomrule
\end{longtable}}
